\section{Objetivos de Negocio}
\label{sec:BusinessObjectives}

Tras la \textbf{captura de requisitos} descrita en la Sección~\ref{sec:RequirementsCaptureMethodology}, la implantación del sistema se orienta a una serie de objetivos de negocio vinculados a la mejora de la actividad diaria del distribuidor y a la relación operativa con las peluquerías clientes. Estos objetivos expresan el \textbf{impacto esperado sobre la forma de trabajar del negocio} una vez la solución se encuentre en uso.

Los principales objetivos de negocio identificados son los siguientes:

\begin{itemize}
    \item \textbf{Reducir la dependencia de procesos manuales} y de \textit{herramientas dispersas} en la gestión comercial, logística y administrativa del distribuidor.
    \item \textbf{Centralizar la información relevante del negocio} para disminuir duplicidades, inconsistencias y pérdida de contexto entre visitas, pedidos, repartos y cobros.
    \item \textbf{Mejorar la eficiencia operativa} del distribuidor, acortando el tiempo dedicado a tareas de coordinación, consulta y seguimiento.
    \item \textbf{Facilitar el control de la actividad comercial diaria} mediante una visión más estructurada de clientes, visitas, rutas y resultados.
    \item \textbf{Aumentar la trazabilidad} de las operaciones críticas del negocio, especialmente en acceso al sistema, gestión de usuarios, seguimiento de pedidos y registro de cobros.
    \item \textbf{Mejorar la calidad del servicio} ofrecido a las peluquerías clientes, permitiéndoles consultar catálogo, pedidos e información de reparto desde un \textit{canal digital unificado}.
    \item \textbf{Favorecer la fidelización} de los clientes profesionales mediante una relación más continua, más transparente y apoyada en herramientas digitales de valor operativo.
    \item \textbf{Sentar una base organizativa y técnica} que permita escalar la actividad comercial e incorporar nuevos usuarios, nuevos comerciales y nuevas funcionalidades sin depender de procedimientos informales.
\end{itemize}

En consecuencia, el valor de negocio esperado no se limita a digitalizar tareas ya existentes, sino a reordenar el \textbf{flujo completo de información del distribuidor} para dotarlo de mayor control, trazabilidad y capacidad de ampliación.
